% !TEX root = cubic_first_order_deformations.tex
\documentclass[a4paper,11pt]{ltjsarticle}
\usepackage[margin=25mm,headheight=16pt]{geometry}
\usepackage{amsmath,amssymb,amsthm,mathtools}
\usepackage{booktabs,enumitem}
\usepackage[unicode,hidelinks]{hyperref}
\usepackage{fancyhdr}
\usepackage{microtype}
\setlength{\emergencystretch}{2em}
\allowdisplaybreaks[2]
\setlist{itemsep=0.2em,topsep=0.4em}
\pagestyle{fancy}\fancyhf{}
\fancyhead[L]{\small 中心三次対の一次変形と二次障害}
\fancyhead[R]{\small 2026-11-13 / v1}
\fancyfoot[C]{\thepage}
\newtheorem{theorem}{定理}[section]
\newtheorem{proposition}[theorem]{命題}
\newtheorem{lemma}[theorem]{補題}
\newtheorem{corollary}[theorem]{系}
\theoremstyle{definition}
\newtheorem{definition}[theorem]{定義}
\newtheorem{example}[theorem]{例}
\theoremstyle{remark}
\newtheorem{remark}[theorem]{注意}
\renewcommand{\proofname}{証明}
\numberwithin{equation}{section}

\hypersetup{pdftitle={中心三次対の一次変形とPlücker二次障害}}
\title{中心三次対の一次変形とPl\"ucker二次障害}
\author{}
\date{2026年11月13日／研究草稿 v1}
\begin{document}
\maketitle
\begin{abstract}
滑らかな三次超曲面を定める多項式環と中心三次元の対について、
双対数環上の変形を調べる。一次の非可換積を曲率付きHochschild境界として消去し、
内部次数シフトと三次周期性を保つmatrix factorizationの比較を構成する。
五変数では、結合積の二次延長障害の係数イデアルがPl\"ucker二次式のイデアルと一致する。
さらに双対数基底上の線束残余圏との比較を付録に収め、
既知のinfinitesimal categorical Torelliから圏論的接空間への完全列を得る。
一次変形、二次障害の零スキーム、全高次の形式変形空間を区別する。
\end{abstract}

\section*{問題設定と既知結果との関係}
多項式環の結合積と指定された中心三次元を同時に変形する。
剰余の多項式環との同定は固定し、同型はその同定の下で恒等になるものに限る。
主な対象は次数別有限自由な双対数変形であり、
一般の有限パラメータの非可換三次超曲面を分類するものではない。

一次変形の議論はHochschild--Kostant--Rosenberg同型を用いるが、
Kontsevichの全次数formalityそのものは必要としない\cite{Kontsevich}。
曲率付きbar写像の行列因子分解への拡張は
C\u{a}ld\u{a}raru--Tuのtwisted complexの記述\cite{CT}に沿って有限式で行う。
滑らかな三次三様体のPoisson構造とGrassmannianとの対応は既知であり、
Mayanskiyの定理と係数計算\cite[Theorem~1, Section~2]{PMayanskiy}と比較する。
ここで扱うのは、その二次式が中心対の結合積の延長障害として現れることの直接的な説明である。
圏論的接空間の同型はJacovskis--Lin--Liu--Zhangの既知定理\cite[Theorem~4.4(2)]{JLLZ}を用いる。
これらの組合せの初出・新規性の包括的な確認は未完了である。

一般の変数数$N\ge5$について一次消去と接空間を扱い、
Pl\"ucker表示および圏論的接空間の完全列は$N=5$に限定する。
後者の証明が使う相対stabilization、Homの底変更、生成性、rotationの比較は
付録\ref{fdw:sec:window}で証明する。
非正則な双対数基底で、全coherent圏の絶対特異圏へ通常のOrlov定理を適用することはない。

% First-order cancellation, quadratic obstructions, and categorical tangent map.
\section{一次の非可換変形と可換模型の接空間}
\label{fpk:sec:first-order}

この節では一般点の全次数形式比較を前提とせず、
双対数環上で非可換積の一次摂動が曲率付きMFにどう作用するかを直接調べる。
これは全次数の定理\ref{cft:thm:complete}の一次部分の独立検証でもある。
曲率の一次変形を固定しなければ、任意の非可換変形が自明という意味ではない。

$S=\mathbb C[x_1,\ldots,x_N]$、$N\ge5$、
$D=\mathbb C[\varepsilon]/(\varepsilon^2)$ とする。
$F\in S_3$は滑らかな射影三次超曲面を定めるものとする。
$A_D$を単位付き・次数別有限自由な$D$-代数で
$A_D/\varepsilon A_D\simeq S$、
$\Omega_D\in Z(A_D)_3$の剰余が$F$となるものとする。

\begin{theorem}[一次Poisson方向の曲率付き消滅]
\label{fpk:thm:infinitesimal}
上の入力に対し、単位と内部次数を保つ$D$-加群標識を適切に選べば、
$A_D$の積と中心元は
\begin{equation}\label{fpk:eq:normal}
 a*b=ab+\varepsilon\,\Theta_\pi(a,b),\qquad
 \Omega_D=F+\varepsilon G,\qquad
 \Theta_\pi(a,b)=\frac12\{a,b\}_\pi
\end{equation}
と書ける。ここで$\pi$は二次係数のbivectorで$[\pi,F]=0$である。
さらに、この標識で有限rank graded MFの
$D$-線形unital $A_\infty$同値
\begin{equation}\label{fpk:eq:equivalence}
 \operatorname{MF}^{\mathbb Z}(A_D,\Omega_D)
 \simeq
 \operatorname{MF}^{\mathbb Z}(S\otimes D,F+\varepsilon G)
\end{equation}
が存在し、剰余では恒等、内部shiftおよび
中心元による三次周期同型$(3)\simeq[2]$を保つ。

二つの標識で得た$G$の差は
$J(F)_3=(\partial_1F,\ldots,\partial_NF)_3$に属する。
従って、一次の可換模型の座標同値類を定める写像は
\begin{equation}\label{fpk:eq:tangent}
 [G]\in(S/J(F))_3
\end{equation}
を通る。特に$G=0$と取れる正規化では、一次Poisson変形は
rotation付きgraded MF圏の一次変形を生じない。
\end{theorem}

\begin{proof}
次数別自由な加群標識を選ぶと、積は$\mu+\varepsilon\beta$となる。
単位を保つよう標識を取り、$\beta$をnormalizedな内部weight零の
Hochschild二cochainとする。結合律の一次部分は
$d_{\mathrm{Hoch}}\beta=0$である。
polynomial algebraのHKR同型により、そのクラスは
二次係数bivector$\pi$の反対称微分cochain$\Theta_\pi$で表せる。
差$\beta-\Theta_\pi$は内部weight零のHochschild coboundaryなので、
$\mathrm{id}+\varepsilon T$という一次の厳密な加群標識変更で
積を式\eqref{fpk:eq:normal}にできる。
この変更は元の代数の結合律を保ち、単位も保つ。
中心元の表示を$F+\varepsilon G$と書く。
中心性の一次部分は$[\pi,F]=0$である。

$F$が滑らかなので、全偏微分は正則列である。
$\iota_{dF}$による斉次Koszul複体をrank二・係数次数二で用いると、
$[\pi,F]=0$から定数係数三重vector
$\Xi\in\bigwedge^3\langle\partial_1,\ldots,\partial_N\rangle$を得て
\begin{equation}\label{fpk:eq:primitive}
 \pi=[\Xi,F]
\end{equation}
となる。原像の差を消すにはrank四・係数次数$-2$の原像が必要となり、
そのような多項式はないので$\Xi$は一意である。
ここでは$\pi$のJacobi恒等式を使わない。

$\Theta_\Xi=\operatorname{HKR}(\Xi)$を、標準の$1/3!$反対称化で
正規化した三入力Hochschild cochainとする。
三つの座標方向の単純三重vectorについては
\[
 \Theta_\Xi(a,b,c)
 =\frac16\det
 \begin{pmatrix}
  \partial_i a&\partial_i b&\partial_i c\\
  \partial_j a&\partial_j b&\partial_j c\\
  \partial_k a&\partial_k b&\partial_k c
 \end{pmatrix}.
\]
式\eqref{fpk:eq:primitive}と反対称性より
\begin{equation}\label{fpk:eq:insert}
 \Theta_\Xi(F,a,b)-\Theta_\Xi(a,F,b)+\Theta_\Xi(a,b,F)
 =\Theta_\pi(a,b).
\end{equation}
$\Theta_\Xi$は多重導分なので
$d_{\mathrm{Hoch}}\Theta_\Xi=0$である。

周期変数に$\deg_{\mathrm{coh}}t=2$、
$\deg_{\mathrm{int}}t=-3$を与え、曲率を$-tF$とする。
\[
 \eta=\varepsilon\,t^{-1}\Theta_\Xi
\]
はcohomological・内部次数ともに零、arity三のnormalized cochainである。
曲率付きGerstenhaber括弧で式\eqref{fpk:eq:insert}を使うと
\[
 [\eta,\mu]=0,\qquad
 [\eta,-tF]=-\varepsilon\Theta_\pi
\]
となる（符号は式\eqref{fpk:eq:insert}を採用する）。
従って、曲率付き構造
$\mu-tF+\varepsilon(\Theta_\pi-tG)$は、
$\eta$による一次ゲージで$\mu-t(F+\varepsilon G)$へ移る。
$\varepsilon^2=0$だから、そのbar写像は
$F_0=0,F_1=\mathrm{id},F_3=\varepsilon t^{-1}\Theta_\Xi$のみを持ち、
逆は$F_3$の符号を反転させればよい。

有限rank MFの行列$Q$は$\varepsilon$の倍数である必要はない。
しかし対象像$Q'=Q+F_3(Q,Q,Q)$および射のtwisting挿入公式は
双対数環上では有限和である。
bar恒等式がMF方程式と関手恒等式を与え、逆bar写像で互いに逆な
$A_\infty$関手が得られる。写像は内部weight零で
$t$を微分しないので内部shiftと$(3)\simeq[2]$を保つ。
これで式\eqref{fpk:eq:equivalence}が従う。

最後に、積を$\mu+\varepsilon\Theta_\pi$へ正規化する
二つの標識の差を$\mathrm{id}+\varepsilon T$とする。
同じ正規化された積を保つ条件は
$d_{\mathrm{Hoch}}T=0$であるから、$T$は内部weight零の導分である。
従って$T(x_i)=\sum_j a_{ij}x_j$であり、
中心元の一次係数は$G\mapsto G+T(F)$と変わる。
この差は$J(F)_3$に属する。
また、同じ$F$の一次座標変更による変形は
この形の$T(F)$で尽くされる。
従って式\eqref{fpk:eq:tangent}を得る。
\end{proof}

\begin{corollary}[Jacobian接空間の次元]
\label{fpk:cor:dimension}
滑らかな三次式$F$に対して
\begin{equation}\label{fpk:eq:dimension}
 \dim_{\mathbb C}(S/J(F))_3
 =\binom{N+2}{3}-N^2.
\end{equation}
特にcubic threefoldでは$10$、cubic fourfoldでは$20$である。
\end{corollary}
\begin{proof}
偏微分の二次式は正則列であり、一次係数のsyzygyを持たない。
従って$J(F)_3$は$N$個の二次偏微分に$N$次元の一次式を掛けた
$N^2$次元の部分空間である。
一方$\dim S_3=\binom{N+2}{3}$である。
\end{proof}

\begin{remark}[Jacobi条件が必要になる位置]
\label{fpk:rem:jacobi}
一次の曲率付き同値に必要なのは中心性からの$[\pi,F]=0$であり、
$[\pi,\pi]=0$は必要ない。
二次以降で同じゲージを延長するには
$[\Xi,\pi]=0$が必要になる。
滑らかな三次式の下では、Jacobi条件$[\pi,\pi]=0$から
勾配Koszulと係数次数により$[\Xi,\pi]=0$を導ける
（補題\ref{cft:lem:primitive}）。
従って一次消滅と全次数の比較は異なる論理的強さを持つ。
\end{remark}


\subsection{一次対変形の接空間と二次のPl\"ucker障害}

\begin{proposition}[次数付き中心三次対の一次接空間]
\label{fpk:prop:tangent-pair}
$N\ge5$とし、可換対$(S,F)$における次数別flatな中心三次対の
双対数変形を、剰余で恒等となる単位付き次数付き同型で割る。
中心元は次数三の実際の元として標識する。
この一次変形集合には自然なベクトル空間構造があり、
HKR正規形を選べば
\begin{equation}\label{fpk:eq:tangent-pair}
 T^1_{\mathrm{pair},(S,F)}
 \simeq
 \bigwedge^3\mathbb C^N\oplus\operatorname{Jac}(F)_3
\end{equation}
と表される。第一成分は$\pi=\iota_{dF}\Xi$、第二成分は
$\Omega=F+\varepsilon G$の$[G]$である。
特に$N=5$ではこの接空間は$20$次元である。
\end{proposition}
\begin{proof}
定理\ref{fpk:thm:infinitesimal}で用いたnormalized HKRにより、
積の一次項の同型類は二次係数bivector$\pi$で記述される。
中心性の一次条件は$[\pi,F]=0$であり、勾配Koszulの完全性と
係数次数から唯一の定数三重vector$\Xi$が存在して$\pi=[\Xi,F]$となる。
逆に任意の$\Xi$はこの中心性条件を満たし、
$\Theta_\pi$を一次積とする双対数上の結合代数を与える。
任意の三次式$G$は中心元の一次項に取れる。
一方、積のHKR正規形を保つ剰余恒等の次数付き標識変更は
一次の導分$T$であり、$G$を$G+T(F)$へ移す。
内部次数零の導分は一次形式の線形変換だから、
$T(F)$の張る空間は$J(F)_3$である。
よって表示を得る。
\end{proof}

以下では$N=5$とする。固定体積形式を使い、
$\bigwedge^3V\simeq\bigwedge^2V^*\otimes\det V$
で$\Xi$に対応する二形式を$\alpha_\Xi$と書く。
その五つのPfaffianを$q_0,\ldots,q_4$とし、
\[
 q(\Xi)=\alpha_\Xi\wedge\alpha_\Xi
 \in\bigwedge^4V^*
\]
で表す。$q(\Xi)=0$は$\Xi$の分解可能性と同値で、
$q_i$は$\operatorname{Gr}(3,V)$のPl\"uckerイデアルを生成する。

\begin{theorem}[二次Jacobi障害のスキーム論的Pl\"ucker表示]
\label{fpk:thm:plucker-obstruction}
$V=\mathbb C^5$、$F\in\operatorname{Sym}^3V^*$を
滑らかな射影三次超曲面の方程式とする。
定数三重vector$\Xi$について$\pi_\Xi=\iota_{dF}\Xi$と置くと
\[
 [\pi_\Xi,\pi_\Xi]=0
 \quad\Longleftrightarrow\quad
 \Xi\text{は分解可能（零を含む）}.
\]
さらに$[\pi_\Xi,\pi_\Xi]$の全係数を$\Xi$の座標の二次式と見て
生成されるイデアルは、五つのPl\"ucker二次式のイデアル
$(q_0,\ldots,q_4)$に\emph{スキームとして}一致する。
\end{theorem}
\begin{proof}
まず$\Xi$が分解可能なら、座標変換により
$\Xi=\partial_0\wedge\partial_1\wedge\partial_2$とできる。
$\pi_\Xi$は三つのactive変数のJacobian bracketであり、
任意の$F$についてJacobi恒等式を満たす。

逆に$\Xi$が分解不可能なら、五次元における二形式の交代標準形より、
座標変換と非零スカラー倍で
\[
 \Xi=\partial_0\wedge\partial_1\wedge\partial_2+
       \partial_0\wedge\partial_3\wedge\partial_4
\]
とできる。$F_i=\partial_iF$、$F_{0i}=\partial_0\partial_iF$
と書く。$\pi_{ij}=\sum_k\Xi_{ijk}F_k$を用いてJacobiatorを直接展開すると、
$[\pi_\Xi,\pi_\Xi]=0$はちょうど
\begin{equation}\label{fpk:eq:jacobi-normal}
 F_{0i}F_j-F_{0j}F_i=0
 \qquad(0\le i<j\le4)
\end{equation}
に等しい（成分の順序による符号と共通非零係数は無関係）。
式\eqref{fpk:eq:jacobi-normal}に$x_j$を掛けて$j$について加え、
Euler等式$\sum_jx_jF_j=3F$および
$\sum_jx_jF_{0j}=2F_0$を使えば
\begin{equation}\label{fpk:eq:euler-obstruction}
 3F F_{0i}=2F_0F_i
 \qquad(0\le i\le4)
\end{equation}
を得る。もし$F_0=0$なら$F$は$x_0$に依らず、
$[1:0:\cdots:0]$が特異点となる。
$F_0\ne0$なら、滑らかな射影三次式$F$は既約なので
$\gcd(F,F_0)=1$である。式\eqref{fpk:eq:euler-obstruction}から
$F\mid F_i$が全$i$について従うが、$\deg F_i=2<3$なので
全$F_i=0$となり矛盾する。従って$\Xi$は分解可能である。

次に二次式の\emph{イデアル}の一致を示す。
$[\pi_\Xi,\pi_\Xi]$の係数は$\Xi$について斉次二次式であり、
分解可能な$\Xi$で零になる。Grassmannianの斉次イデアルは
五つのPl\"ucker二次式で生成されるので、
係数全体は$q_0,\ldots,q_4$の線形結合である。
従ってある線形写像$L_F$について、Jacobiatorの全係数ベクトルは
$L_F(q(\Xi))$と表される。

ここで二次写像$q:\bigwedge^3V\to\bigwedge^4V^*$は全射である。
実際、五次元では任意の非零四形式は分解可能であり、
適当な基底で$c\,e^1\wedge e^2\wedge e^3\wedge e^4$と書ける。
$\alpha=e^1\wedge e^2+(c/2)e^3\wedge e^4$とすれば
$\alpha\wedge\alpha$はその四形式である。
もし$L_F$に非零の核があれば、$q(\Xi)$がその非零元となる$\Xi$を選べる。
すると$[\pi_\Xi,\pi_\Xi]=0$なのに$q(\Xi)\ne0$であり、
前半の結論と矛盾する。よって$L_F$は単射である。
その双対写像の全射性から、Jacobiatorの係数二次式は
Pl\"ucker二次式の五次元空間全体を張り、生成イデアルも一致する。
\end{proof}

\begin{corollary}[二次まで延長できる非可換一次方向]
\label{fpk:cor:second-order}
$N=5$の中心三次対の一次変形の非可換方向
$\Xi\in\bigwedge^3V$が、次数付き\emph{結合積}の二次変形
$\mathbb C[\varepsilon]/(\varepsilon^3)$へ延長可能であるためには、
五つのPl\"ucker二次式が零でなければならない。
結合積に関する二次障害の零スキームはPl\"ucker錐と一致する。
その錐の局所次元は$7$で、頂点のZariski接空間は$10$次元である。
これは中心元を含めた対の高次延長障害がすべて消えるという主張ではない。
\end{corollary}
\begin{proof}
結合積$\mu+\varepsilon\beta_1+\varepsilon^2\beta_2$の結合律の
二次項は
$d_{\rm Hoch}\beta_2+\frac12[\beta_1,\beta_1]=0$である。
HKR同型で二次障害の$HH^3$類は
$\frac12[\pi_\Xi,\pi_\Xi]$に対応し、
この類が零であることが結合積の二次延長の必要十分条件である。
前定理からその係数イデアルは五つのPl\"ucker関係式になる。
$\operatorname{Gr}(3,5)$の次元は$3(5-3)=6$なので
錐の次元は$7$、零点では全五式の一次微分が零で接空間は$10$次元である。
\end{proof}

\begin{remark}[Pl\"ucker障害の射程]
\label{fpk:rem:plucker-scope}
滑らかな三次三様体のPoisson構造とGrassmannianの対応は
既知の研究\cite{PMayanskiy}と整合する。
上の定理は、その対応を\emph{結合積の二次延長障害の係数イデアル}として
直接記述するものである。
この記述の先行研究における初出は未確認である。
障害の零スキームを同定したことだけから、
中心元も含む対の全高次形式変形関手を同定したとは結論しない。
\end{remark}


\subsection{Kuznetsov成分の圏論的接空間への射}
\begin{theorem}[可換点での圏論的忘却射の微分]
\label{fpk:thm:categorical-tangent}
$N=5$、$X=V(F)\subset\mathbb P^4$を滑らかなcubic threefoldとする。
次数別flatな中心三次対の一次変形を、剰余で恒等な次数付き同型で割った
$T^1_{\mathrm{pair},(S,F)}$を考える。
Artin基底上のenhancedなstabilization/window比較
（定理\ref{am:thm:comparison}）によって、
rotation付き残余圏の変形に付随する
\emph{基礎dg圏を忘れた}Kodaira--Spencer写像
\[
 \kappa_{\rm cat}:T^1_{\mathrm{pair},(S,F)}
 \longrightarrow HH^2(\operatorname{Ku}(X))
\]
を定めることができる。
命題\ref{fpk:prop:tangent-pair}のHKR正規形による分解では
\begin{equation}\label{fpk:eq:categorical-KS}
 \kappa_{\rm cat}(\Xi,[G])=\eta([G]),
 \qquad
 \eta:H^1(X,T_X)\xrightarrow{\ \sim\ }HH^2(\operatorname{Ku}(X)).
\end{equation}
特に次の完全列を得る。
\begin{equation}\label{fpk:eq:categorical-exact}
 \boxed{
 0\longrightarrow\bigwedge^3\mathbb C^5
 \longrightarrow T^1_{\mathrm{pair},(S,F)}
 \xrightarrow{\ \kappa_{\rm cat}\ }
 HH^2(\operatorname{Ku}(X))
 \longrightarrow0 .
 }
\end{equation}
したがって、可換点での圏論的忘却射の微分は全射で、
その核は$10$次元、像も$10$次元である。
この結論は非零Poisson一般点での微分の階数を主張しない。
\end{theorem}
\begin{proof}
滑らかな三次三様体の標準正規束完全列とEuler列によって、
三次式の変形の接空間は
\[
 H^1(X,T_X)\simeq
 S_3/\langle x_i\partial_jF:0\le i,j\le4\rangle
 =\operatorname{Jac}(F)_3
\]
と同定される。一次Poisson方向は命題\ref{fpk:prop:tangent-pair}の
$\Xi$で表される。定理\ref{fpk:thm:infinitesimal}は、
同じ双対数環上で、剰余が恒等でrotationと指定周期を保つ
MFの実際の$A_\infty$同値を
$(A_D,\Omega_D)$と$(S_D,F+\varepsilon G)$の間に構成する。
Artin基底上のenhanced windowと可換側の相対Orlov比較を合成すると、
この同値は両者の残余dg圏の一次変形を同一視する。
従って、その$HH^2$におけるKodaira--Spencer類は
非可換Poisson方向$\Xi$には依存せず、
滑らかな三次式$F+\varepsilon G$のKuznetsov成分の変形類
$\eta([G])$に等しい。

Jacovskis--Lin--Liu--Zhang
\cite[Theorem~4.4(2)]{JLLZ}は、任意の滑らかな
index-two・degree-three Fano threefoldについて
$\eta:H^1(T_X)\to HH^2(\operatorname{Ku}(X))$が同型であることを示す。
従って式\eqref{fpk:eq:categorical-KS}から
式\eqref{fpk:eq:categorical-exact}が従う。
ここで$HH^2$は\emph{基礎dg圏}の変形接空間であり、
rotationを忘れずに扱うモジュライstack全体の接空間と
自動的に同一視したものではない。
\end{proof}
\begin{remark}[圏論的微分と相対復元]
\label{fpk:rem:categorical-scope}
完全列\eqref{fpk:eq:categorical-exact}は\emph{可換点での}
基礎dg圏の変形接空間に関するものである。
非零Poisson一般点に沿う忘却射の微分を記述するには、
相対的な可換復元、中心元の標識とKodaira--Spencer写像の比較を
さらに必要とする。これは可換点の計算とは別の問題である。
\end{remark}

\begin{remark}[モジュライの局所次元と接空間]
\label{fpk:rem:moduli}
滑らかな三次三様体では$\dim\operatorname{Jac}(F)_3=10$である。
これに対し、\emph{非零Poisson一般点の次元}と
\emph{可換点でのZariski接空間}は異なる。
固定した滑らかな五変数三次式に対するPoisson錐はPl\"ucker錐であり、
非零分解可能な一般点では$1+\dim\operatorname{Gr}(2,5)=7$次元である。
一方、錐の頂点では全てのPl\"ucker方程式が二次式なので、
Zariski接空間は$\bigwedge^3\mathbb C^5$の全$10$次元となる。
元の可換cubicの変形接空間$10$次元を合わせると、
中心三次対の可換点における一次接空間は$20$次元である。
二次障害零スキームの局所次元$17$はこの接空間の次元と異なる。
本節の二次障害計算はその差を直接説明するが、
対の\emph{全高次}形式局所環の同定をこの計算だけから結論しない。

定理\ref{fpk:thm:categorical-tangent}の完全列は可換点での核が
$10$次元であることを示すが、非零Poisson一般点での核の次元は与えない。
後者については注意\ref{fpk:rem:categorical-scope}を参照されたい。
\end{remark}

\appendix
\section{双対数基底上の線束残余圏とMFの比較}
\label{fdw:sec:window}

この節は圏論的接空間の結論に必要な相対比較を収録する。
$D=\mathbb C[\varepsilon]/(\varepsilon^2)$、$d\ge5$とし、
$A_0=D$、各$A_j$が有限自由で、$A/\varepsilon A=\mathbb C[x_1,\ldots,x_d]$となる
次数付き中央$D$-代数$A$を考える。
$\Omega\in Z(A)_3$の剰余$F$は滑らかな射影三次超曲面を定め、
$B=A/(\Omega)$、$a=d-3$、$i=1-a$と置く。
各dg圏は有限coneと直和因子について閉じたものとする。
全coherent圏ではなく線束で生成される部分を扱い、
非正則な$D$上で絶対特異圏への比較を適用しない。

\subsection{次数別自由性と相対stabilization}

\begin{lemma}[相対Koszul分解と中心商]\label{fdw:lem:resolution}
$A,B$は左右noetherianで、$D$は左右の次数付き$A$-加群として
$P_j=A(-j)^{\binom dj}$、$0\le j\le d$という有限自由分解を持つ。
各内部次数でこの分解は$D$-分裂し、剰余への底変更後も完全である。
$\Omega$は正則で、$B_j$は$D$-自由である。
\end{lemma}
\begin{proof}
$\varepsilon$による有限filtrationの随伴次数環は
$\mathbb C[x_1,\ldots,x_d]\otimes D$なので、有限filtrationによる
持上げから左右noetherian性を得る。
$A\to D$のaugmentation核は次数別$D$-自由である。
その剰余の生成元を持ち上げて$A(-1)^d$から全射を作る。
次数別の余核が$D$-自由だから核も$D$-自由で、剰余へのtensorは完全性を保つ。
これを繰り返すと、剰余が多項式環のKoszul syzygyであることから、
$j$段目には次数$j$の$\binom dj$個の生成元だけが必要になる。
$d$段の後の核の剰余は零なので、次数別Nakayamaによって核は零である。
同じ議論は左加群にも通る。
$\Omega:A_{j-3}\to A_j$の剰余は非零多項式による単射である。
有限自由$D$-加群の行列の単元小行列式を用いれば分裂単射となり、
各次数の余核$B_j$が自由であることと$\Omega$の正則性が従う。
ここから$A$の有限大域次元を結論してはいない。
\end{proof}

$\operatorname{Res}_A$を制限関手とし、標準dg enhancementで
\[
 \mathcal D_{B/A}=
 \{M\in D^b(\operatorname{gr}B):\operatorname{Res}_AM\in\operatorname{Perf}(\operatorname{gr}A)\},
 \qquad
 \mathcal S_{B/A}=(\mathcal D_{B/A}/\operatorname{Perf}(\operatorname{gr}B))^{\rm idem}
\]
と定める。$B$は$A$上長さ一の分解を持つので、このdg商は定義される。

\begin{lemma}[enhancedな余核比較]\label{fdw:lem:stabilization}
有限自由graded MFの余核関手は、内部shiftと三次周期を保つdg Morita同値
$\operatorname{MF}^{\mathbb Z}(A,\Omega)\simeq\mathcal S_{B/A}$を与える。
その逆を$\operatorname{stab}$と書く。比較は$D\to\mathbb C$への底変更と両立する。
\end{lemma}
\begin{proof}
有限次数付き射影$A$-加群は、そのaugmentation商の自由基底を持ち上げ、
次数付きNakayamaを使えば有限次数付き自由である。
$\mathcal E$を、有限生成次数付きGorenstein-projective $B$-加群のうち
$A$上有限射影次元を持つもののexact圏とする。
Chenのchange-of-rings定理\cite[Theorem~2.10, Corollary~2.11]{AMChen}から、
その非零対象の$A$-射影次元は一である。
syzygyとcosyzygyでこの性質が保たれるので、$\mathcal E$は
有限射影$B$-加群をprojective-injectiveとするFrobenius圏である。

有限自由MFの成分別分裂exact圏を$\mathcal F$とする。
余核は$\mathcal F\to\mathcal E$のexact関手である。
逆に$M\in\mathcal E$について
$0\to P_1\xrightarrow{\varphi}P_0\to M\to0$を取ると、
$\Omega M=0$より$\Omega\operatorname{id}_{P_0}$は$\varphi$を通る。
$\varphi$の単射性により他方の積も$\Omega\operatorname{id}_{P_1}$となり、MFを得る。
射の持上げとhomotopyを含めた安定圏の同値は
\cite[Theorem~7.2]{AMChen}による。斉次分解を選べば次数付きでも同じ証明が通る。

enhancementは三角同値だけから宣言しない。
上のexact関手が誘導するdg商の関手
\[
 D^b_{\rm dg}(\mathcal F)/K^b_{\rm dg}(\operatorname{proj\!-\!inj}\mathcal F)
 \longrightarrow
 D^b_{\rm dg}(\mathcal E)/K^b_{\rm dg}(\operatorname{proj}B)
\]
を使う。Frobenius分解により$H^0$は安定圏であり、全shiftでの同値から
全Homのcohomologyも同型になる。
左のdg商とMFの比較は有限複体のtotalizationによる。
成分別splitなacyclic複体で、横微分とcontractionを$d,s$、
縦のMF微分を$\delta$とし、$ds+sd=1$、$s^2=0$、
$d\delta+\delta d=0$、$\delta^2=w$とする。
横方向が有限なので$q=s(1+\delta s)^{-1}=(1+s\delta)^{-1}s$は有限和であり、
\[
 (1+s\delta)\{1-(d+\delta)q-q(d+\delta)\}(1+\delta s)
 =-sws=0.
\]
よってtotalizationは可縮で、dg商を通る。

右のdg商から$\mathcal S_{B/A}$への関手は加群圏への包含による。
完全射影分解から全shiftのstable Homが特異圏のHomと一致する。
$\mathcal D_{B/A}$の任意の対象は高い$B$-syzygyのshiftと特異圏で同型である。
このsyzygyは$A$上有限射影次元を持ち、Chenの定理によって$B$上有限
Gorenstein-projective次元を持つため、さらにsyzygyを取れば$\mathcal E$に入る。
これで本質的全射性を得る。

有限自由MFの余核の各次数について、剰余上のfactorization mapは単射である。
実際、その二つの積は非零多項式$F$である。
単元小行列式から余核は次数別$D$-自由となり、
余核・exact列・有限totalizationは剰余への底変更と両立する。
以上のdg関手のzigzagとMorita逆により、所要の自然な比較を得る。
内部shiftと指定中心元に由来する周期も各段階で保存される。
\end{proof}

\subsection{線束Hom、proper性、底変更}

$\pi:\operatorname{Gr}B\to\operatorname{QGr}B$、
$\mathcal O(j)=\pi B(j)$、導来sectionを$\mathbf R\omega$とする。

\begin{lemma}[線束の有限window]\label{fdw:lem:lines}
線束の導来Homは$D$-perfectで剰余への底変更と両立し、
\[
 \operatorname{Ext}^q(\mathcal O(r),\mathcal O(s))=
 \begin{cases}B_{s-r}&q=0,\\B_{r-s-a}^{\vee}&q=d-2,\\0&\text{その他}\end{cases}
\]
となる。ここで$(-)^\vee=\operatorname{Hom}_D(-,D)$、負の添字は零とする。
$0\le j\le d-1$について
$(\mathbf R\omega\mathcal O(j))_{\ge i}=B(j)_{\ge i}$である。
$E=\bigoplus_{j=0}^{d-1}\mathcal O(j)$は線束生成部分
$\mathcal C_B=\operatorname{thick}_D\{\mathcal O(j):j\in\mathbb Z\}$を生成する。
\end{lemma}
\begin{proof}
補題\ref{fdw:lem:resolution}の有限Koszul分解を双対化する。
固定内部次数では有界有限自由$D$-複体で、剰余のKoszul双対性を
単元成分の消去で持ち上げると
\[
 \mathbf R\underline{\operatorname{Hom}}_A(D,A)=D(d)[-d],\qquad
 \mathbf R\underline{\operatorname{Hom}}_B(D,B)=D(a)[-(d-1)]
\]
を得る。後者には正則中心元のchange-of-ringsを使う。
$\chi$も確認する。$B/(\varepsilon)$の$B$-自由分解は
$\varepsilon$の乗法による周期分解であり、有限生成$M$のHom cohomologyは
有限生成$B/(\varepsilon)$-加群である。
change-of-rings spectral sequenceと可換剰余環の$\chi$、
$D$の有限な剰余体filtrationを用いれば、
$\operatorname{Ext}_B^q(D,M)$は各$q$で内部次数について有限である。
従ってArtin--Zhangのsection--torsion三角を適用できる\cite{AZ}。

$B/B_{\ge\ell}$の次数filtrationの商は$B_j\otimes_DD(-j)$である。
上の双対性と$B_j$の自由性より、固定内部次数$t$では$j=-t-a$の一層だけが寄与し、
$\ell>j$で安定する。従って
\[
 H^q(\mathbf R\tau B)_t=
 \begin{cases}B_{-t-a}^{\vee}&q=d-1,\\0&\text{その他}\end{cases}
\]
となる。三角$\mathbf R\tau B\to B\to\mathbf R\omega\pi B$からHom公式を得る。
$\mathbf R\omega\mathcal O(j)$の正のcohomologyは$t\le-j-a$にしかなく、
$i=1-a>-j-a$なので指定切断後に消える。

底変更には、次数別$D$-自由なsyzygyを持つ$B/B_{\ge\ell}$の分解を用いる。
固定次数のHom複体の剰余はcohomologyが一つの次数にしかない。
剰余上のsplit contractionを$\varepsilon$進の有限摂動で持ち上げれば、
有限自由cohomologyの一項複体へのhomotopy同値を得る。
従って通常のtensorが導来tensorを計算する。
Hom、filtered colimit、section三角を通じて自然な底変更写像が準同型となる。
filtered colimitはK-flat性を保つので、無限複体の平坦性を無条件に仮定していない。

Koszul分解を$B$へ導来tensorした複体のcohomologyは$D$と$D(-3)$のshiftであり、
いずれも射影商で消える。よって長さ$d$の線束複体は完全である。
両端の階数一を用いれば、連続した$d$本から全てのtwistを前後に生成できる。
Hom公式を有限coneと直和因子へ延長すれば、$\mathcal C_B$のproper性と底変更も従う。
\end{proof}

\begin{lemma}[MF Homの有限縮約と生成子の持上げ]\label{fdw:lem:mf}
有限自由graded MFのHomは$D$-perfectで、剰余への導来底変更と両立する。
$P$の剰余が古典graded MFを生成するなら、$P$も$D$上で全MFを生成する。
\end{lemma}
\begin{proof}
周期変数$t$の二重次数を$(2,-3)$とする。
有限rank行列のHomのcohomological次数と内部次数を固定すると$t$の冪が定まり、
その成分は有限個の$A_j$の直和である。
剰余Homは滑らかな三次超曲面のgraded MFであるから、
古典Orlov比較によって有界有限次元cohomologyを持つ\cite{Orlov}。
この剰余複体のcohomologyへのsplit contractionを選ぶ。
微分の差を$\delta\in\varepsilon\operatorname{End}$とすれば、
$(1+s\delta)^{-1}$と$(1+\delta s)^{-1}$は二乗零性により有限和であり、
homological perturbationによって有界有限自由複体へ縮約できる。
これはK-flat性と自然なHom底変更を与える。

生成性は右直交の消滅だけで済ませず、$\mathcal D=\operatorname{Ind}(\operatorname{MF})$で示す。
Homの底変更と剰余$\bar P$の生成性により、カテゴリーの剰余は古典MFとMorita同値である。
$P$はcompactであり、そのHomはcolimitと可換する。
$M\in\mathcal D$が$P$と右直交なら、
$\operatorname{RHom}(\bar P,M\otimes_D^{\mathbf L}\mathbb C)=0$なので
$M\otimes_D^{\mathbf L}\mathbb C=0$である。
$D$は$D(D)$において$\mathbb C$の二段のextensionだから、$M=0$となる。
よって$P$は$\mathcal D$のcompact generatorであり、そのcompact部分でも生成する。
\end{proof}

\subsection{windowとrotationの比較}

\begin{theorem}[双対数上のenhanced window]\label{am:thm:comparison}
$\mathcal K_B^{\rm lin}=\mathcal C_B\cap
\langle\mathcal O,\ldots,\mathcal O(a-1)\rangle^\perp$とする。
関手$\Psi=\operatorname{stab}\circ(\mathbf R\omega-)_{\ge i}$はdg Morita同値
\[
 \Psi:\mathcal K_B^{\rm lin}\xrightarrow{\sim}\operatorname{MF}^{\mathbb Z}(A,\Omega)
\]
を与える。これは剰余への底変更と両立し、
$\rho=L_{\mathcal O}\circ(1)$を内部shiftへ移す。
指定中心元の周期性は$\rho^3\simeq[2]$へ移る。
\end{theorem}
\begin{proof}
線束Hom公式から$\mathcal O,\ldots,\mathcal O(a-1)$は$D$-例外列であり、
evaluation coneによる射影$p$を持つ。従って$G=p(E)$が残余圏を生成する。
derived sectionと切断、補題\ref{fdw:lem:stabilization}のMorita逆から$\Psi$を構成する。
$0\le j\le d-1$について、その入力は$B(j)_{\ge i}$である。
短完全列$0\to B(j)_{\ge i}\to B(j)\to T_j\to0$の
有限長商$T_j=B(j)_{<i}$は、有限自由$D$の内部shiftから構成される。
$D$と$B$は$A$上perfectだから、三項とも$\mathcal D_{B/A}$に属する。
従って$\Psi$は有限coneと直和因子を通して$\mathcal C_B$全体に定義される。

$j<a$では$T_j=0$なので、$\Psi$は例外部分を消す。
$j=a,a+1,a+2$では$\Psi(\mathcal O(j))=\operatorname{stab}(T_j)[-1]$で、
\[
 \operatorname{gr}T_j=\bigoplus_{r=0}^{j-a}B_r\otimes_DD(j-r).
\]
$B_0=D$より、三つの像は三角的に三つのstabilizationを生成する。
剰余では$\Psi$は古典Orlovのsection--windowであり、
例外部分によるVerdier商からMFへの同値になる\cite[Theorem~6.4]{BS}。
左直交と右直交は、同じVerdier商を表示することを介して比較する。

sourceのHomは補題\ref{fdw:lem:lines}、targetのHomは補題\ref{fdw:lem:mf}により
$D$-perfectで底変更と両立する。自然なHom比較のconeの剰余は零なので、
導来Nakayamaによりcone自体も零である。これで完全忠実性を得る。
$\Psi(G)$の剰余は古典MFのgeneratorであり、補題\ref{fdw:lem:mf}によって
$D$上でも全MFを生成するので、本質的全射性が従う。

$M$が残余圏に属するとき、$j\ge1$の直交性は
$\operatorname{RHom}(\mathcal O(j),M(1))=0$と線束Hom公式から、
$j=0$の直交性は左mutationの定義から従う。よって$\rho$は残余圏を保つ。
また
$(\mathbf R\omega M)_i=\operatorname{RHom}(\mathcal O(a-1),M)=0$
なので、自然な切断三角により
\[
 ((\mathbf R\omega M)(1))_{\ge i}
 =(\mathbf R\omega M)_{\ge i+1}(1)
 \simeq(\mathbf R\omega M)_{\ge i}(1).
\]
$\Psi(\mathcal O)=0$とevaluation三角を併せると
$\Psi\rho\simeq(1)\Psi$を得る。
内部shiftは自己同値なので$\rho$も自己同値であり、周期性はMFの$t$から移る。
以上は全coherent圏の任意の底変更定理を用いていない。
\end{proof}

\begin{thebibliography}{99}
\bibitem{Kontsevich}
M. Kontsevich, \emph{Deformation quantization of Poisson manifolds},
Lett. Math. Phys. \textbf{66} (2003), 157--216.
\url{https://arxiv.org/abs/q-alg/9709040}.
Section~4.6.2のHKR比較を参照する。

\bibitem{CT}
A. C\u{a}ld\u{a}raru and J. Tu,
\emph{Curved $A_\infty$ algebras and Landau--Ginzburg models},
New York J. Math. \textbf{19} (2013), 305--342.
Section~2, Remark~2.8, Theorem~2.12.
\url{https://nyjm.albany.edu/j/2013/19-16v.pdf}.

\bibitem{PMayanskiy}
E. Mayanskiy, \emph{Poisson cohomology of two Fano threefolds},
arXiv:1303.5916, Theorem~1およびSection~2の係数計算。
\url{https://arxiv.org/abs/1303.5916}.
同Theorem~1はLoray--Pereira--Touzetの既知結果として述べられている。

\bibitem{JLLZ}
A. Jacovskis, X. Lin, Z. Liu and S. Zhang,
\emph{Infinitesimal categorical Torelli theorems for Fano threefolds},
J. Pure Appl. Algebra \textbf{227} (2023), 107418.
\url{https://doi.org/10.1016/j.jpaa.2023.107418}.
著者版Theorem~4.4(2), 印刷p.~15。
\url{https://ajacovskis.github.io/hochschild_only_2.pdf}.

\bibitem{AMChen}
X.-W. Chen, \emph{Module factorizations}, arXiv:2402.11613v3 (2025).
Theorem~2.10, Corollary~2.11, Theorem~7.2。
\url{https://arxiv.org/html/2402.11613v3}.
余核の三角同値を使用し、enhancementは本文でexact圏とdg商により構成する。

\bibitem{AZ}
M. Artin and J. J. Zhang, \emph{Noncommutative projective schemes},
Adv. Math. \textbf{109} (1994), 228--287.
\url{https://doi.org/10.1006/aima.1994.1087}.

\bibitem{Orlov}
D. Orlov, \emph{Derived categories of coherent sheaves and triangulated categories of singularities},
arXiv:math/0503632v3 (2009), Theorem~2.5。
\url{https://arxiv.org/abs/math/0503632}.

\bibitem{BS}
J. Burke and G. Stevenson, \emph{The derived category of a graded Gorenstein ring},
arXiv:1507.00830v1 (2015), Theorem~6.4およびその証明。
\url{https://arxiv.org/abs/1507.00830}.
付録では剰余体上のwindowとVerdier商の向きの同定に用いる。

\end{thebibliography}
\end{document}
